
\label{fpinversionsection}The matrix inversion procedure is based on the
possibility of matrix factorization as discussed in Sec.\ref{LUsection}.
Formally, the iterative method may be expressed in the simple form 
\begin{equation}
\widehat{\Bbb{N}}^{\prime }Z^{(k+1)}=\widehat{\Bbb{N}}^{\prime
}Z^{(k)}-\beta _{C}\left( \widehat{\Bbb{N}}^{\prime }Z^{(k)}-\widehat{\Bbb{P}%
}_{N}^{-1}\Upsilon ^{(k)}\right)   \label{fpitequation}
\end{equation}
where $\beta _{C}$ is the iteration parameter that may be adjusted for an
improvement of the rate of convergence, using a Chebyshev acceleration
technique as shown in Ref. \cite{kar86}. Here $\widehat{\Bbb{P}}$ is used
for the main diagonal precondiotiong as explained in Sec.\ref{LUsection} .
However, it has been shown that values of $\beta _{C}$ other than unity do
not give much better results in general for implicit methods (\cite{pey98}),
and therefore, only the case $\beta _{C}=1$ is considered. Then Eq. \ref
{fpitequation} is equivalent to Eq. \ref{generalmatrixform} given in Sec. 
\ref{LUsection}. Replacing $\widehat{\Bbb{N}}^{\prime }$ by $\widehat{\Bbb{L}%
}\widehat{\Bbb{U}}$, one obtains 
\begin{equation}
\widehat{\Bbb{L}}\widehat{\Bbb{U}}Z^{(k+1)}=\widehat{\Bbb{P}}%
_{N}^{-1}\Upsilon ^{(k)}  \label{LUfpitequation}
\end{equation}

If we introduce the increment $\Delta _{Z}^{(k+1)}$%
\begin{equation}
\Delta _{Z}^{(k+1)}=Z^{(k+1)}-Z^{(k)}
\end{equation}
and the residual $R_{Z}^{(k)}$%
\begin{equation}
R_{Z}^{(k)}=\widehat{\Bbb{P}}_{N}^{-1}\Upsilon ^{(k)}-\widehat{\Bbb{N}}%
^{\prime }Z^{(k)}
\end{equation}
Eq.\ref{LUfpitequation} becomes simply 
\begin{equation}
\widehat{\Bbb{L}}\widehat{\Bbb{U}}\Delta _{Z}^{(k+1)}=R_{Z}^{(k)}
\end{equation}
which can be solved by two successive inversion steps, evaluating
respectively forward and backward solutions from triangular systems of
equations 
\begin{eqnarray}
\widehat{\Bbb{L}}V_{Z}^{(k+1)} &=&R_{Z}^{(k)} \\
\widehat{\Bbb{U}}\Delta _{Z}^{(k+1)} &=&V_{Z}^{(k+1)}
\end{eqnarray}

For achieving convergence towards the steady state solution $Z^{(\infty )}$,
various iterative methods may be used, like the Conjugate Gradient Squared 
(CGS), the BiConjigate Gradient (BICG) or the Quasi-Minimal Residual (QMR)
methods. These methods, which both give equivalent results in term of
convergence rate, are preferred to the BiConjigate Gradient Stabilized
(BICGSTAB) method, since the conservative scheme is nearly always well
fullfiled even for marginaly well-conditionned matrices, like in presence of
RF waves whose intensity is large. All methods here considered are build-in
MatLab functions whose principle may be found in Refs. \cite{pre92} or \cite
{axe94} for more detailed insights.

After numerous tests, it has been found that the drop tolerance parameter $%
\delta _{lu}$ has a close link with the maximum number of iterations $MAXIT$
allowed inside the function that performs the iterative matrix inversion, in
order to avoid cumulative numerical errors that lead usually to violation of
the conservative nature of the code. Even if this problem may be cured by
forcing the Maxwellian solution close to $p=0$ and normalizing the solution
at each iteration, this approach has not been considered, because it may
hinder more serious problems regarding the overall numerical conservative
scheme. So far, a very robust and fully conservative scheme is achieved with
the following rule of thumb
\begin{equation}
\left\{ 
\begin{array}{l}
\delta _{lu}=10^{-5}:MAXIT<2-3 \\ 
\delta _{lu}=10^{-4}:MAXIT<4-5 \\ 
\delta _{lu}=10^{-3}:MAXIT<9-10 \\ 
\delta _{lu}=10^{-2}:MAXIT<19-20
\end{array}
\right. 
\end{equation}
for all cases addressed by the code, including the in presence of RF waves.
More detailed studied are necessary to clarify this point.

At each iteration, $Z^{(k+1)}$ is evaluated from the knowledge of $\Delta
_{Z}^{(k+1)},$ and the electron distribution function $%
f_{0,l+1/2,i+1/2,j+1/2}^{\left( 0\right) }$ is reconstructed in order to
evaluate the non-linear term $\widehat{C}\left( X_{M},Z^{(k)}\right) $ that
arises from self-collisions. Though this procedure is quite time consuming,
its incidence on the overall time for converging is fairly marginal, since
the number of iterations never exceeds $10$ in most cases.

According to Ref. \cite{biz97}, convergence towards $Z^{(\infty )}$ is
considered to be achieved and the iteration is stopped when the following
criterion is fullfiled 
\begin{equation}
\sqrt{\frac{2\pi \iint \left[ \partial f_{0}^{\left( 0\right) }/\partial
t\right] ^{2}f_{0}^{\left( 0\right) }J_{p}J_{\xi _{0}}dpd\xi _{0}}{2\pi
\iint f_{0}^{\left( 0\right) }J_{p}J_{\xi _{0}}dpd\xi _{0}}}<\mathcal{R}_{%
\mathcal{C}}^{f}  \label{fpconvcrit}
\end{equation}
for an arbitrarily small $\mathcal{R}_{\mathcal{C}}^{f},$ at all plasma
locations $\psi ,$ where $J_{p}$ and $J_{\xi _{0}}$ are the partial
Jacobians as introduced in Sec. \ref{fluxcons}\footnote{%
The possibility to define a fully $3-D$ criterion has been foreseen.
However, since collision rates may change drasticaly from plasma center to
the edge, this choice ensures that convergence is well achieved at the right
level at all plasma locations.}.

Projected on the numerical grid, the criterion (\ref{fpconvcrit}) becomes

\begin{equation}
\sqrt{\frac{\sum_{i=0}^{i=n_{p}-1}\sum_{j=0}^{j=n_{\xi _{0}}-1}\left[
f_{0,l+1/2,i+1/2,j+1/2}^{\left( 0\right) \left( k+1\right)
}-f_{0,l+1/2,i+1/2,j+1/2}^{\left( 0\right) \left( k\right) }\right]
^{2}f_{0,l+1/2,i+1/2,j+1/2}^{\left( 0\right) \left( k\right)
}p_{i+1/2}^{2}\Delta p_{i+1/2}\Delta \xi _{0,j+1/2}}{\Delta
t^{2}\sum_{i=0}^{i=n_{p}-1}\sum_{j=0}^{j=n_{\xi
_{0}}-12}f_{0,l+1/2,i+1/2,j+1/2}^{\left( 0\right) \left( k\right)
}p_{i+1/2}^{2}\Delta p_{i+1/2}\Delta \xi _{0,j+1/2}}}<\mathcal{R}_{\mathcal{C%
}}^{f}
\end{equation}

The criterion introduced in Ref. \cite{kar86} has the main advantage to
ensure that convergence is achieved not for electrons whose energy is close
to the thermal one, but also for the fastest when a tail is driven by an
external source. Usually, $\mathcal{R}_{\mathcal{C}}^{f}$ ranges between $%
10^{-10}$ and $10^{-12},$ and in the code the standard value is $10^{-11}.$
The quantity (\ref{fpconvcrit}) may be viewed as a norm, and its definition
is consequently independent of the level of current which results from the
the lack of symmetry of $f_{0}^{\left( 0\right) }$ in momentum space. It may
therefore be used for strong deformations of the electron distribution
function, even if the net current is close to zero.

In practice, the total number of iteration $\mathcal{N}_{\mathcal{C}}^{f}$
is a free choice. It is usually set to $50$, so that inversion procedure
stops when (\ref{fpconvcrit}) is fullfiled in quite all cases. If the number
of iterations reaches $\mathcal{N}_{\mathcal{C}}^{f},$ the program warns
that convergence is not achieved and that something is going wrong in the
calculations. This may occur sometimes when too large quasilinear RF
diffusion coefficient are used, leading to inconsistent and non physical
solutions. This point is extensively discussed in Sec.\ref{lhcdsection}.
Conversely, when time step is very large, i.e. $\Delta \overline{t}\geq 1000,
$ the convergence may be achieved for a given value $\mathcal{R}_{\mathcal{C}%
}^{f}$  before the current level is fully established. Such an event may
arise because there is a too reduced number of iterations which involves the 
$1^{st}$ order Legendre non-linear correction term for momentum
conservation. Consequently, to avoid this problem, a minimum number of
iterations has been set to $6$, which turns out to be enough for giving
reliable results, owing to the weakness of the non-linearity.
